% Options for packages loaded elsewhere
\PassOptionsToPackage{unicode}{hyperref}
\PassOptionsToPackage{hyphens}{url}
%
\documentclass[
]{article}
\usepackage{lmodern}
\usepackage{amssymb,amsmath}
\usepackage{ifxetex,ifluatex}
\ifnum 0\ifxetex 1\fi\ifluatex 1\fi=0 % if pdftex
  \usepackage[T1]{fontenc}
  \usepackage[utf8]{inputenc}
  \usepackage{textcomp} % provide euro and other symbols
\else % if luatex or xetex
  \usepackage{unicode-math}
  \defaultfontfeatures{Scale=MatchLowercase}
  \defaultfontfeatures[\rmfamily]{Ligatures=TeX,Scale=1}
\fi
% Use upquote if available, for straight quotes in verbatim environments
\IfFileExists{upquote.sty}{\usepackage{upquote}}{}
\IfFileExists{microtype.sty}{% use microtype if available
  \usepackage[]{microtype}
  \UseMicrotypeSet[protrusion]{basicmath} % disable protrusion for tt fonts
}{}
\makeatletter
\@ifundefined{KOMAClassName}{% if non-KOMA class
  \IfFileExists{parskip.sty}{%
    \usepackage{parskip}
  }{% else
    \setlength{\parindent}{0pt}
    \setlength{\parskip}{6pt plus 2pt minus 1pt}}
}{% if KOMA class
  \KOMAoptions{parskip=half}}
\makeatother
\usepackage{xcolor}
\IfFileExists{xurl.sty}{\usepackage{xurl}}{} % add URL line breaks if available
\IfFileExists{bookmark.sty}{\usepackage{bookmark}}{\usepackage{hyperref}}
\hypersetup{
  hidelinks,
  pdfcreator={LaTeX via pandoc}}
\urlstyle{same} % disable monospaced font for URLs
\setlength{\emergencystretch}{3em} % prevent overfull lines
\providecommand{\tightlist}{%
  \setlength{\itemsep}{0pt}\setlength{\parskip}{0pt}}
\setcounter{secnumdepth}{-\maxdimen} % remove section numbering

\date{}

\begin{document}

(relativity4)=

\hypertarget{lesson-18-brspecial-relativity-4}{%
\section{\texorpdfstring{Lesson 18: \\
Special Relativity,
4}{Lesson 18:  Special Relativity, 4}}\label{lesson-18-brspecial-relativity-4}}

\begin{quote}
Einstein's Mass and Energy
\end{quote}

\begin{verbatim}
Here we go. The T-shirt equation. Nothing is the same after this fourth of Einstein's 1905 papers...that one that said $m=\dfrac{E}{c^2}$. We'll develop it without calculus (although if you want to see a full derivation, it's included for those with that background). Energy now simplifies! Mass becomes interesting! Particle collisions become understandable! Matter gets made! And I get my white board on the Big Bang Theory!
\end{verbatim}

\hypertarget{a-little-bit-of-minkowski}{%
\subsection{A Little Bit of Minkowski}\label{a-little-bit-of-minkowski}}

\begin{verbatim}
---
align: left
figclass: margin
---
 Hermann Minkowski <br> 1864-1909<br>
\end{verbatim}

\begin{quote}
"It came as a tremendous surprise, for in his student days Einstein had
been a lazy dog... He never bothered about mathematics at all."
\end{quote}

The story of Hermann Minkowski is both engaging and sad. He grew up in
Königsberg, Germany after some time in Russia, to a middle-class
business family. His mathematical skills were apparent in high school
and he entered the University of Königsberg where he stayed through his
PhD degree, winning a prestigious prize at the age of 17 for a solution
to a challenge problem issued in European mathematics. Mathematicians
like to do that sort of thing.

In 1887 he was recruited to the University of Bonn where he taught for
seven years and returned to his alma mater, the University of
Königsberg, where he stayed for two years before he was appointed the
chair of mathematical physics at Eidgenössische Polytechnikum Zürich in
1894. He married in 1897 and had two daughters, the second in 1902.

\begin{verbatim}
---
align: left
figclass: margin
---
 Hermann Minkowski during the last year of his life.
\end{verbatim}

That was a big year for him professionally as well as David Hilbert
created a chair for Minkowski at the center of European mathematics at
the University of Göttingen. It was there that he developed the ideas of
space and time that outlived him and set the path for a firm foundation
of Special Relativity and motivated the approach to the General Theory
of Relativity. This mathematization of physics was not originally
welcome to Einstein, but he grew to embrace it. In fact, in some ways it
was Minkowski's style and his skill that changed the relation of physics
and mathematics to the present day.

\begin{quote}
Morris Kine in a review wrote, "\emph{A key point of the {[}Minkowski{]}
paper is the difference in approach to physical problems taken by
mathematical physicists as opposed to theoretical physicists. In a paper
published in} 1908 \emph{Minkowski reformulated Einstein's 1905 paper by
introducing the four-dimensional (}space-time*) non-Euclidean geometry,
a step which Einstein did not think much of at the time. But more
important is the attitude or philosophy that Minkowski, Hilbert, Klein,
and Weyl pursued, namely, that purely mathematical considerations,
including harmony and elegance of ideas, should dominate in embracing
new physical facts. \underline{Mathematics so to speak was to be master
and physical theory could be made to bow to the master.} {[}my
emphasis{]} Put otherwise, theoretical physics was a subdomain of
mathematical physics, which in turn was a sub-discipline of pure
mathematics."
\end{quote}

That Minkowski embraced the implications of Special Relativity is a
delicious irony as he was Einstein's mathematics professor in Zurich and
his opinion of him was right in line with that of the other faculty:

\begin{quote}
"Oh, that Einstein, always cutting lectures... I really would not
believe him capable of it."
\end{quote}

It was Minkowski who merged space and time into a single
"\emph{Spacetime}" continuum -- a mathematical space, but importantly, a
model of how the universe is really stitched together. That's the
"Spacetime" in \emph{Quarks, Spacetime, and the Big Bang}. Below we'll
visit some of "Minkowski Space" and here's a preview from his notes.
Remember the plot at the top.

\begin{verbatim}
---
width: 80%
figclass: margin-caption
alt: My figure text
name: minkowskispace
---
MARGINCAP
\end{verbatim}

This work he did in 1908. While it was well-received, it was a
semi-popular lecture that he gave that was the stimulus to many to take
Relativity seriously. Ask Mr Google about "views of space and time" and
he'll reward you with 3.5 billion hits:

\begin{quote}
"The views of space and time which I wish to lay before you have sprung
from the soil of experimental physics, and therein lies their strength.
They are radical. Henceforth, space by itself, and time by itself, are
doomed to fade away into mere shadows, and only a kind of union of the
two will preserve an independent reality."
\end{quote}

To the shock of the community and dismay to his many friends around the
world, Minkowski died suddenly at the age of 44 due to a ruptured
appendix.

\hypertarget{working-our-way-toward-relativistic-energy}{%
\subsection{Working Our Way toward Relativistic
Energy}\label{working-our-way-toward-relativistic-energy}}

At the end of the last lesson, we crept up on the notion that the speed
of an object being pushed at a constant value changes, but not linearly,
so that it never exceeds the speed of light. Here's a cartoonish sketch
of that summarizes what seems to be going on.

What we had is a constant force pushing on a mass, creating a momentum,
and so a speed.

\begin{verbatim}
---
width: 70%
figclass: margin-caption
alt: My figure text
name: block_limited
---
MARGINCAP
\end{verbatim}

Now, let's compare the Newtonian experience (rather, assumptions) with
the relativistic experience.

\begin{verbatim}
---
width: 70%
figclass: margin-caption
alt: My figure text
name: limited
---
MARGINCAP
\end{verbatim}

In (a) we have the Newtonian notion of a fixed mass and an unlimited
speed and so an unlimited momentum. In (b) we have the relativistic
situation in which we know that the speed is limited and there's no
limitation on the momentum. So how does that happen? The mass must be
unlimited. Not constant, as in the Newtonian definition.

\begin{verbatim}
---
width: 60%
figclass: margin-caption
alt: My figure text
name: screenshot_2224
---
MARGINCAP
\end{verbatim}

Mass is going to become a very strange thing.

\begin{quote}
Remember that mass for Newton was problematic. It was a definition that
was circularly defined in terms of density. How you measure it presumes
that you know how to measure acceleration and force...without the mass.
Circular. Relativity is the opening salvo against the naive notion of
mass from Newton and we'll see that with the Higgs Boson discovery in
2012, our ideas of mass changed again...to be even more strange.
\end{quote}

\hypertarget{relativistic-momentum}{%
\subsubsection{Relativistic Momentum}\label{relativistic-momentum}}

Momentum is still:

\begin{verbatim}
:class: warning

$$
p = m\times \frac{\text{a change in distance}}{\text{a change in time}} = m \frac{\Delta x_H}{\Delta t_H} \nonumber
$$

But by now you know the song: that numerator and denominator depend on the frame and so a new kind of momentum – relativistic momentum – is required. I'll just state it here:

$$
p = m\gamma u
$$

Looks familiar, but there's the $\gamma$ function, which tends towards $1$ when $u$ is very small, so we get back "regular" momentum.
\end{verbatim}

\hypertarget{an-einstein-thought-experiment}{%
\subsubsection{An Einstein Thought
Experiment}\label{an-einstein-thought-experiment}}

Let's follow one of Einstein's typically imaginative thought experiments
that he suggested in his original 1905 fourth paper. The T-shirt
equation paper. In order to appreciate it, I need to tell you something
that he postulated in his first 1905 paper that was on light quanta --
that light consists of little particles that he called "quanta" and
which were named "photons" in 1925 by someone else. The thing that we
need from this idea (which we will talk about a lot later) is that these
little particles carry momentum.

So inside a box is a source of light (photons) at the left and an
absorber of photons on the right. We start with the box sitting still,
minding its own business.

\begin{verbatim}
---
width: 90%
figclass: margin-caption
alt: My figure text
name: box1
---
MARGINCAP
\end{verbatim}

Now the bulb is turned on and light begins to stream away from it, that
is, photons carry momentum to the right.

\begin{verbatim}
---
width: 90%
figclass: margin-caption
alt: My figure text
name: box2
---
MARGINCAP
\end{verbatim}

So momentum is now imbalanced to the right, unless the box itself moves
to the left with the same, but oppositely directed value of momentum. It
recoils.

So from the outside, you see this box suddenly begin to lurch to the
left! Then, when the photons strike the absorber on the right-hand side
of the box (perfectly absorbing so that there's no reflection)...the box
suddenly stops!

\begin{verbatim}
---
width: 90%
figclass: margin-caption
alt: My figure text
name: box3
---
MARGINCAP
\end{verbatim}

Einstein analyzed this event and showed that the center of mass of the
box rearranged itself and that therefore the photon has to behave as if
it has a mass.

Seriously. In the first paper (and every other paper anyone has written
since!) the photon is an object with zero mass. And yet, this simple
head-experiment concludes that the light has an energy -- even Maxwell
knew about electromagnetic waves carrying energy. But since light has
momentum, we must interpret the situation to suggest that the photon's
energy leads to a second conclusion that suggest that energy is
mass-like. That's a lot to swallow.

\hypertarget{approximating-functions}{%
\subsubsection{Approximating Functions}\label{approximating-functions}}

There are a number of ways to get to the T-shirt equation and most
require calculus which I'll not do in \emph{QS\&BB}. I want to sneak up
on it with a plausibility argument, rather than a direct proof. In
\textbackslash@ref(\#lessonmath) I introduced the notion of
approximating a function's shape by successively adding together a
string of functions, each one getting you closer and closer to the full
function. This was Newton's idea. What I did in that lesson was look at
the function

\[f(x) = \frac{1}{1-x} \nonumber\]

and showed that it can be represented as the series

\[f(x) = \frac{1}{1-x} = 1+x+x^2 + x^3 +... \nonumber\]

the dot-dot-dot means that this progression continues to an infinite
number of terms and that when you add them all up you get back
\(\dfrac{1}{1-x}\). Almost magical.

Here's what that looks like:

\begin{verbatim}
---
width: 90%
figclass: margin-caption
alt: My figure text
name: binom_ex
---
MARGINCAP
\end{verbatim}

Notice a couple of things. First, the full function is the top, bold red
curve. Second, notice that each curve in the series does indeed get
closer and closer to the desired function. It's going to take a lot of
them to get there -- an infinite number of them.

:::\{admonition\} Pencils Out! 🖋 📓 :class: warning

Finally, notice that for small \(x\), say below \(x=0.2\) the first two
terms, \(1+x\) do a pretty good job. So in practical terms, if you're
sure that your \(x\)'s are going to be less than that value you don't
have to use the whole function, you can get by with just \(1+x\).

So, with that reminder...we have a function that's in every relativistic
equation, namely, \(\gamma\):

\[\gamma = \frac{1}{\sqrt{1-\beta^2}} \nonumber\]

and maybe you know where I'm going. We can indeed express this function
in an infinite series and it would be:

\[\gamma = \frac{1}{\sqrt{1-\beta^2}} = 1+\frac{\beta^2}{2}+\frac{3\beta^4}{8}+\frac{5\beta^6}{16}+...\]

Here it is plotted:

\begin{verbatim}
---
width: 90%
figclass: margin-caption
alt: My figure text
name: binom_gamma
---
MARGINCAP
\end{verbatim}

where I've emphasized only the first two terms. Again, if we are only
interested in very low \(\gamma\)'s we can get away with just the first
couple of terms and we've done a pretty good job of approximating the
whole relativistic gamma function. Knowing that "regular life" is pretty
Newtonian, that part of the gamma function should do a pretty good job
for barely relativistic situations.

:::

\hypertarget{einsteins-energy-equation-through-the-back-door}{%
\subsubsection{Einstein's Energy Equation, Through the Back
Door}\label{einsteins-energy-equation-through-the-back-door}}

Let's play with that term now and see if we can gain some insight.

:::\{admonition\} Pencils Out! 🖋 📓 :class: warning

This is a pencil-out activity. Look at the first two terms in the
\(\gamma\) expansion:

\begin{align*}
\gamma \approx& 1 + \beta^2/2 \\
\gamma =& 1 + \frac{1}{2}\frac{u^2}{c^2}
\end{align*}

Look at the second term. This is really suspiciously like the classical
kinetic energy form, \(K=\frac{1}{2}mv^2\), right? Let's run with this
and see what happens. In order to nudge it along, let's multiply by
\(m\) and by \(c^2\). That would do it, right?

\begin{align*}
\gamma =& 1+\frac{1}{2}\frac{u^2}{c^2} \\
mc^2\gamma =& mc^2 + mc^2\frac{1}{2}\frac{u^2}{c^2} \\
m\gamma c^2 =& mc^2 + \frac{1}{2}mu^2
\end{align*}

By now we're used to the idea that slow objects' modeling are an
approximation to the complete relativistic objects' form. We're
obviously in a slow-ish format here by keeping only the first two terms
of the \(\gamma\) function. So we isolated the "regular" kinetic energy
piece on the right and so we also know that \emph{whatever is left in
the other two terms are also energy terms}:

\begin{itemize}
\item
  \(m\gamma c^2\) is some kind of energy
\item
  \(mc^2\) is another kind of energy
\item
  \(\frac{1}{2}mu^2\) is the familiar kind of energy of motion, kinetic
  energy
\end{itemize}

This is a fundamental change in how we think about energy. What we've
created through the back door is a new kind of energy:

\begin{itemize}
\item
  Total energy of an object \(=\) its energy of mass \(+\) its energy of
  motion
\end{itemize}

Even if the object were not moving, \emph{it would still have energy in
this picture}. In fact, all of energy can be described by this little
formula - potential energy, chemical energy, nuclear energy,
electromagnetic radiation energy, the whole shebang energy.

Here's our accounting:

\begin{itemize}
\item
  \(E_T =\) total energy of an object \(= m\gamma c^2\), first
  formulation
\item
  \(E_m =\) mass energy of an object = \(mc^2\)
\end{itemize}

The T-shirt equation:

\begin{verbatim}
---
width: 90%
figclass: margin-caption
alt: My figure text
name: t-shirt
---
MARGINCAP
\end{verbatim}

I've left out the kinetic energy here because my back door approach to
this was for non-relativistic motion, so our original
\(\frac{1}{2}mu^2\) term was relevant. But the fully relativistic
kinetic energy is different and we can simply derive its form from what
we've got by generally calling it "\(K\)."

\begin{align*}
E_T =& E_m + K \\
K =& E_T - E_m \\
K =& m\gamma c^2 - mc^2 \\
K =& mc^2(\gamma -1)
\end{align*}

So our inventory is:

\begin{itemize}
\item
  \(E_T =\) total energy of an object \(= m\gamma c^2\), first
  formulation
\item
  \(E_m =\) mass energy of an object = \(mc^2\)
\item
  \(K =\) energy of motion \(=\) \(mc^2(\gamma -1)\)
\item
  \(p = m\gamma u\) \(=\) relativistic momentum
\end{itemize}

\begin{verbatim}
:class: danger
<a href="../examples/relativity/FILE.html" target="_blank">Derivation of energy relation</a>
\end{verbatim}

With a little bit of algebra, we can write our second form of the total
energy:

\[E_T^2 = (mc^2)^2 + (pc)^2\]

and we can understand Einstein's little game of the box with the light.

:::

\begin{verbatim}
:class: danger

<a href="../examples/relativity/Total_E_1_R4_E1.html" target="_blank">derivation of the total energy</a>
\end{verbatim}

Remember, he had to conclude that if light carries momentum, which
Maxwell knew, then it has to behave as if it has a mass, which is not
strictly true. In Equation \{eq\}\texttt{totalE} we see how that works.
If light quanta are massless, and they are, the total energy is not
zero, it's:

\begin{align*}
E_T^2(\text{for anything}) &= (mc^2)^2 + (pc)^2 \text{ but for a massless photon where } m=0: \\
E_T(\text{photon}) &= pc
\end{align*}

We can see this in a visual way using Pythagoras' theorem since

\[E_T^2 = (mc^2)^2 + (pc)^2  \nonumber\]

\begin{verbatim}
---
width: 60%
figclass: margin-caption
alt: My figure text
name: energytriangle
---
MARGINCAP
\end{verbatim}

For different kinds of objects, as the mass leg gets smaller and
smaller, the hypothenus gets closer and closer to the horizontal leg.

\begin{verbatim}
So what does one do when visiting a friend at UCLA who was the Science Advisor for the television program, The Big Bang Theory? And he hands you a marker and your own on-air whiteboard? Why you tell the world how to expand the gamma function. That's what.
\end{verbatim}

\begin{verbatim}
---
width: 90%
figclass: margin-caption
alt: My figure text
name: bbt
---
MARGINCAP
\end{verbatim}

\begin{verbatim}
:class: danger

<a href="../examples/relativity/Nimitz_penny_1_R4_E3.html" target="_blank">Penny for your thoughts</a>
\end{verbatim}

\hypertarget{two-things-to-worry-about}{%
\subsubsection{Two Things to Worry
About}\label{two-things-to-worry-about}}

There are two issues buried inside of the general total energy relation
that merit attention:

\[E_T^2 = (mc^2)^2 + (pc)^2 \nonumber\]

\begin{verbatim}
**Thing 1.** Remember that Einstein started his discussion with his quantum idea, which we'll have to wait to learn more about. But I've already told you that the photon – the quantum of light – has no mass. Notice that this equation for energy accommodates that.

$$
\begin{align*}
E_T^2 \text{ for } m=& 0 \text{ becomes } (pc)^2 \text{ so }\\
E_T =& pc
\end{align*}
$$

which is precisely what he suggested was the case for a single quantum of light in his first 1905 paper.
\end{verbatim}

and

\begin{verbatim}

**Thing 2.** Remember that a square can have two roots, one positive and one negative. What about the total energy squared?

$$
\begin{align*}
E_T^2 =& (mc^2)^2 + (pc)^2 \text{ means that the roots can be } \\
E_T =& \pm \sqrt{(mc^2)^2 + (pc)^2}
\end{align*}
$$

What are we to make of a negative energy? Stay tuned. But this plagued efforts to create a quantum mechanics that was relativistic.
\end{verbatim}

\begin{verbatim}
:class: danger

<a href="../examples/relativity/Estimate_gamma_1_R4_E2.html" target="_blank">How good is the approximation for gamma?</a>
\end{verbatim}

\hypertarget{remember-how-we-got-here}{%
\subsubsection{Remember How We Got
Here?}\label{remember-how-we-got-here}}

I started worrying about inertia in the last lesson when we found that a
constant force applied to a massive object would result in a constant
acceleration (our example was \(g\)) but that the speed of that object
would increase at a slower and slower rate, never passing \(c\). Our
notion of relativistic mass is a way to think about why that is the
case. Since here, \(m_R = \gamma m\) it increases with speed through the
relativistic gamma function:

Remember to push Run:

So if you're fortunate enough to be moving at 80\% of the speed of
light, your mass as viewed from a Home frame would be about 1.75 times
your rest mass as you would "feel" in the Away frame. Food for thought,
if you're a particle physicist. Stay tuned.

Particle accelerators? They operate in the high-end of gamma...here's a
useful example of the gamma region of interest:

Remember to push Run:

\begin{verbatim}
:class: danger

<a href="https://loncapa.msu.edu//tiny/msu/gmgg9M"  target="_blank">Precision practice with the plots</a>
\end{verbatim}

\hypertarget{what-is-mass-anyway}{%
\subsubsection{What Is Mass, Anyway?}\label{what-is-mass-anyway}}

Einstein's fourth, less than three page 1905 paper in November, was
entitled, \emph{Does The Inertia Of A Body Depend Upon Its Energy
Content?} and he worked out (with different names for the variables):

\begin{quote}
"If a body gives off the energy E in the form of radiation, its mass
diminishes by \(m=\dfrac{E}{c^2}\)." And he closed with, "It is not
impossible that with bodies whose energy-content is variable to a high
degree (e.g. with radium salts) the theory may be successfully put to
the test. If the theory corresponds to the facts, radiation conveys
inertia between the emitting and absorbing bodies."
\end{quote}

This. Is an extraordinary idea. Any object with a mass has a special
energy solely associated with that mass. "Mass-energy" is a common term
for this and captures the notion succinctly. But "energy-mass" is also
applicable as we saw for the massless photon in the box! It had energy,
but that energy conferred on the photon features that dynamically played
the role of a mass.

It's often said -- INCORRECTLY! -- that mass can be converted into
energy. Mass IS energy and energy IS mass. We can convert the units, but
not the actual substance. The best analogy I can think of is currency.

Today, \$1 is convertible into 0.83 EURO. If I were to go into a shop in
the Schiphol Airport in Amsterdam and purchase a magazine I could pay at
the counter in dollars or I could pay in euros -- both are readily
acceptable in that airport. Their values are identical and indeed,
colloquially we can convert dollars into euros and we would see visually
the two currencies are represented by two different kinds of paper and
coins. The \emph{economic value} however is identical and the passage of
one unit into the other is accomplished by a single conversion factor:
0.83. If I have 10 EURO, I can convert that into dollars:

\[X \text{ USD} = (8.3 \times 10^{-1}) \times 10 \text{ EURO}= 8.3 \text{ USD.}\nonumber\]

The conversion of mass and energy units functions the same way. The
\emph{value} of mass-energy is identical, and the conversion of units is
accomplished by a single conversion factor as well: from \(E=mc^2\), the
conversion is

\[c^2 = (3 \times 10^8)^2=9 \times 10^{16}. \nonumber\]

(in the MKS system). A very big number! If I have a 10 kg object, I can
convert that into Joules:

\[x \text{ J} = (9 \times 10^{16})\times 10 \text{ kg} = 9 \times 10^{17} \text{ J.} \nonumber\]

The size of this is hard to wrap your head around. The amount of energy
trapped in even tiny masses is enormous.

\begin{verbatim}
:class: warning

But there's more to this than just the identification of energy with mass and *vise versa*. There's also the interpretation of mass as not constant. From a silly little rearrangement from above:

$$
\begin{align*}
E_T =& E_m + K = mc^2 + K\\
K =& = mc^2(\gamma -1) = m\gamma c^2 - mc^2 \\
E_T =& mc^2 + m\gamma c^2 - mc^2 = m\gamma c^2
\end{align*}
$$

So, another way of looking at the total energy of an object is to combine its $E_m = mc^2$ with $K$ to get:

$$
E_T = m\gamma c^2 \nonumber
$$

We've now got two "kinds" of mass: what we'll call the "relativistic mass" and the "rest mass." The latter is the mass of an object in its own rest frame, like proper time is the time that you read from the watch you're wearing in your own rest frame. That we'll call:

$$
\text{Rest mass } = m = E_m/c^2 \nonumber
$$

The former is what's sometimes called the Relativistic mass:

$$
\text{Relativistic mass } = m_R = \gamma m \nonumber
$$

The Relativistic mass changes as the speed of a frame changes with respect to a Home observer. By a factor of $\gamma$.

In our Home and Away notation, the masses are related as:

$$
\text{Relativistic mass } = m_H = \gamma m_A \nonumber
$$
\end{verbatim}

\begin{verbatim}
This is a controversial point. While there is no difference between two interpretations of mass, professional physicists do not use "relativistic mass" preferring to keep the idea of mass a constant quantity regardless of motion and deal exclusively with the total energy, $E_T$. I think that for our purposes, the relativistic mass is a very nice interpretation which carries a lot of insight along with it. So don't tell my friends, okay?
\end{verbatim}

\hypertarget{momentum-in-relativity}{%
\paragraph{Momentum In Relativity}\label{momentum-in-relativity}}

Remember that I introduced momentum for relativity as if it's a
definition. That's not strictly the case, and again, it's a calculus
derivation so it's not appropriate for this text. But we can circle back
and retain our original idea of momentum as

\[p=mv \nonumber\]

by using the relativistic mass in this relation: \(m_R = \gamma m\)
where \(m\) is the rest mass. Now we get the originally postulated
relativistic momentum:

\[p=m_Rv = m\gamma v \nonumber\]

Remember, that conversion factor is enormous. We need an apple.

\begin{verbatim}
---
width: 60%
figclass: margin-caption
alt: My figure text
name: applemass
---
MARGINCAP
\end{verbatim}

\begin{verbatim}
:class: warning

Remember that in this beloved example, our table is 1 meter above the ground and the apple has a mass of 0.1 kg. The kinetic energy when it hits the ground is 1 J if we approximate $g \approx 10$ m/s$^2$. What is the energy trapped inside of the apple's mass in comparison?

$$
\begin{align*}
E_m =& mc^2 \\
E_m =& (0.1)(3 \times 10^8)^2 \text{ J } = (0.1)(9 \times 10^{16}) \text{ J } \\
E_m =& 9 \times 10^{15} \text{ J }
\end{align*}
$$

...essentially $10^{16}$ times its motion energy as it hits the floor.
\end{verbatim}

Let's go to an old-time particle accelerator -- one in which protons are
accelerated to only 95\% of the speed of light.

\begin{verbatim}
:class: danger

<a href="../examples/relativity/Spark_1_R4_E2.html" target="_blank">electrons from a big spark</a>
\end{verbatim}

\begin{verbatim}
:class: danger

<a href="https://loncapa.msu.edu//tiny/msu/cMSPC"  target="_blank">increase of mass by protons in an accelerator</a>
\end{verbatim}

\begin{verbatim}
:class: danger

<a href="https://loncapa.msu.edu//tiny/msu/ld0WG3"  target="_blank">Airplanes again</a>
\end{verbatim}

\begin{quote}
From this point on, when I refer to "mass," I'll mean the relativistic
mass, \(m_R\). If I mean the mass of an object inside of its own rest
frame, I'll explicitly refer to the "rest mass," \(m\).
\end{quote}

\hypertarget{the-rubber-meets-the-road-particle-creation-and-potential-energies}{%
\subsection{The Rubber Meets the Road: Particle Creation and Potential
Energies}\label{the-rubber-meets-the-road-particle-creation-and-potential-energies}}

In our lesson on energy we learned to represent energy conservation as

\[E_{T (\text{before})}=K_0 + U_0 = K + U = E_{T (\text{after})}\nonumber\]

Kinetic plus potential energies before...must equal kinetic plus
potential energies after some event. That potential energy term was
always a little squirrelly -- like it was added to make things work.

Let's imagine a football tackle. Player \(1\) runs at player \(2\),
tackles him, and they stick together forming a new, single object
"\(12\)." The masses of the two players are identical and their speeds
are equal.

\begin{verbatim}
---
width: 90%
figclass: margin-caption
alt: My figure text
name: tackleR
---
MARGINCAP
\end{verbatim}

\begin{verbatim}
:class: warning

What would a physics student have calculated in 1904? First, she would have said that momentum was conserved:

$$
\begin{align*}
p_{1,0} - p_{2,0} =& p_{12} = 0 \text{, so} \\
p_{1,0} =& - p_{2,0} \\
mv_1 =& - mv_2 \text{ which, since their masses are identical,} \\
v_1 =& -v_2 \text{ so, call each of them just }v = |v_1|=|v_2|
\end{align*}
$$

And she would have conserved energy:

$$
\begin{align*}
K_{1,0} + K_{2,0} =& K_{12}=0 \\
\frac{1}{2}mv_1^2 + \frac{1}{2}mv_2^2 =& 0 \text{, which, since their masses are identical,} \\
& \text{   and since the $v's$ are identical}\\
v_1^2 + v_2^2 =& 2v^2 = 0
\end{align*}
$$

which doesn't make any sense. At this point her professor would mumble about energy lost in heat, the springiness of the two objects causing compression and then a release of that compression by heating and moving the air, which we'd hear as sound...and so on and so on. That is, this completely inelastic collision would conserve total energy, but we would have written it like this:

$$
\begin{align*}
E_T =& K_{1,0} + K_{2,0} = K_{12} + E(\text{lost}) \\
E_T =& K_{1,0} + K_{2,0} = 0 + E(\text{lost}) \\
E_T =& K_{1,0} + K_{2,0} = E(\text{lost, heat, sound, elastic compression, yadda yadda yadda})
\end{align*}
$$

Now let's suppose it's early 1906 after Einstein's realization, and the collision is a for  *relativistic football tackle*.
\end{verbatim}

\begin{verbatim}
:class: warning

In Special Relativity, we would now write:

$$
\begin{align*}
E_{T (\text{before})} =& E_{m, 1 (\text{before})} + K_{1 (\text{before})} + E_{m, 2 (\text{before})} + K_{2 (\text{before})}\\
=& E_{m, 12 (\text{after})}+K_{12 (\text{after})} = E_{T (\text{after})}
\end{align*}
$$

Total energy is still conserved, but how it arranges itself between energy of mass ($E_m$) and energy of motion ($K$) changes.

Now, momentum is still conserved, so the combined objects, 12, stop dead (no pun intended...but if it's relativistic, they're moving pretty fast!). That means that  $K_{12 (\text{after})}=0$ and we now have:

$$
\begin{align*}
E_{T (\text{before})} =& E_{m, 1 (\text{before})} + K_{1 (\text{before})} + E_{m, 2 (\text{before})} + K_{2 (\text{before})}\\
=& m_1c^2 + K_{1 (\text{before})} + m_2c^2 + K_{2 (\text{before})} = m_{12}c^2
\end{align*}
$$

Let's rearrange this a little:

$$
\begin{align*}
E_{T (\text{before})} =& m_1c^2 + m_2c^2 + K_{1 (\text{before})}  + K_{2 (\text{before})} = m_{12}c^2 = E_{T (\text{after})}
\end{align*}
$$

which is different from the 1904 calculation. Remember, there we found that

$$
K_{1 (\text{before})}  + K_{2 (\text{before})} = 0 \text{ in 1904} \nonumber
$$

and that's now not the situation. In fact,

$$
K_{1 (\text{before})}  + K_{2 (\text{before})} >0 \nonumber
$$

and so

$$
m_{12}c^2  > m_1c^2 + m_2c^2 \text{ ! in 1906} \nonumber
$$
\end{verbatim}

The combination of the two moving objects is more massive when they
collide than they were separately, before they collided. All of that
mumbo jumbo about leftover energy can be interpreted as a mass increase:
the energy of motion has become additional energy of mass. Said, another
way:

\begin{align*}
\text{before 1905, nobody would have disputed that: } &m_{12} = m_1 + m_2 \\
\text{after 1905, we must assert that : } &m_{12} > m_1 + m_2
\end{align*}

\begin{verbatim}
:class: important

Mass, by itself, is not conserved. Mass energy plus energy of motion is always conserved.
\end{verbatim}

\hypertarget{particle-physics-the-sum-is-worth-more-than-the-parts}{%
\subsubsection{Particle Physics: the sum is worth more than the
parts}\label{particle-physics-the-sum-is-worth-more-than-the-parts}}

One of the discoveries in particle physics that I was involved in was
the discovery of the most massive elementary particle we know of: the
"top quark." We'll learn all about that later, but what's relevant here
is that we collided protons with the antiparticles of protons.
Antiparticles are identical to their particle counterpart, except they
have the opposite electrical charge. So they have the same mass.

Remember that I said I'd sometimes refer to the mass of a proton as "1"
in dimensionless units? This will be more of an issue in later lessons,
but some nomenclature:

\begin{verbatim}
* Antiparticles will be written with a "bar" over the top.
* So an electron is $e$ and an anti electron would be $\bar e$. (It's so useful in medicine and other areas of life that it has a name: the positron.)
* A proton is $p$ and so an antiproton is $\bar{p}$.
* A top quark is $t$ and ...you get it, right? The anti-top quark is $\bar t$.
\end{verbatim}

Here's where relativity is useful as a tool. Our top quark discovery
reaction was:

\[p + \bar{p} \to t + \bar{t}. \nonumber\]

Here's a cartoon, where the "before" is at the top, then the physics of
annihilation and creation is some very short time later (middle) and the
two top quarks are the result, at the bottom.

\begin{verbatim}
---
width: 90%
figclass: margin-caption
alt: My figure text
name: ttbarcartoon
---
MARGINCAP
\end{verbatim}

Here's a kind of surprising fact:

\begin{align*}
m(p) = 1 \\
m(\bar p) = 1 \\
m(t) = 173 \\
m(\bar t) = 173
\end{align*}

So if you add up the masses in the initial state of the colliding
proton-antiproton set you get \(m(\text{before})=2\) and the masses of
the resulting top-quark pair, \(m(\text{after}) = 2 \times 173 = 346!\)

That's why we must accelerate the protons and antiprotons to very, very
high energies -- so that the energy of motion in the beams can be turned
into energy of mass in the final products. Sweet.

Here's a colorful model for the collision of two objects into two
objects:

:::\{admonition\} Pencils Out! 🖋 📓 :class: warning

\begin{verbatim}
---
width: 90%
figclass: margin-caption
alt: My figure text
name: twototwo
---
MARGINCAP
\end{verbatim}

:::

Look at each term carefully. Because we're going to create new matter.

Let's make a cartoon particle accelerator, in the spirit of the Large
Hadron Collider in Geneva. Two protons are accelerated around a huge,
evacuated pipe (steered by magnets and accelerated in electric fields as
we've discussed). In this play-accelerator, they are brought together in
one place where they collide.

\begin{enumerate}
\def\labelenumi{\arabic{enumi}.}
\item
  At the top, we see them coming toward one another
\item
  At the bottom they've disappeared and created two, new...things,
  \(t\). The things are much more massive than the original protons.
\end{enumerate}

\begin{verbatim}
:class: warning

**Glad you asked:** Not quite and as maybe you've heard me say before, stay tuned. We will talk about that when we get to quantum mechanics where the explanation is both beautiful and weird.
\end{verbatim}

\begin{verbatim}
---
width: 60%
figclass: margin-caption
alt: My figure text
name: thing12
---
MARGINCAP
\end{verbatim}

Here is the reaction formula for a proton from the Left colliding with a
proton from the Right to make two "things," \(t(1)\) and \(t(2)\):

\[p(L) + p(R) \to t(1) + t(2). \nonumber\]

Feynman Diagram for the process:

\begin{verbatim}
---
width: 60%
figclass: margin-caption
alt: My figure text
name: FD_things
---
MARGINCAP
\end{verbatim}

\begin{verbatim}
:class: warning

**Glad you asked:** That's a quantum mechanics model that we'll get to. Stay tuned. Sorry.
\end{verbatim}

:::\{admonition\} Pencils Out! 🖋 📓 :class: warning

Here are the ground rules for this game. Although energies would
typically be Joules, or electron volts, we'll just use fake energy units
scaled to the mass-energy of a proton.

\[\text{Mass energy of a proton } = E_m = m_pc^2 = 1 \nonumber\]

The mass of a "thing" is big:

\begin{align*}
\text{Let's assume that the mass of 1 thing } =& 3.5\times m_p \\
\text{So, the mass-energy of 1 thing } =& 3.5 \times 1 = 3.5\\
\end{align*}

Nature only produces pairs of "things" so the minimum mass-energy
required to create two things is:

\[\text{Mass energy of two things } = 2 \times 3.5 = 7. \nonumber\]

Finally, the protons are moving -- fast. So there's kinetic energy in
their motions of three times the mass-energy of a proton:

\[\text{Let's assume that the kinetic energy of each proton } = 3 \times m_pc^2 = 3 \nonumber\]

From our energy conservation equation we can keep track of the before
energies and the after energies:

\begin{align*}
E_T(\text{beam protons} =& \text{ Mass energies } + \text{ kinetic energies} \\
E_T =& 2 \times m_pc^2 + 2 \times K \\
E_T =& 2 \times 1 + 2 \times 3 = 2 + 6 = 8
\end{align*}

The question is whether there is enough total energy in the beam protons
to be able to make two "things" in the collision.

You'll need a worksheet that you'll find on-line in the course pack. It
looks like this:

\begin{verbatim}
---
width: 90%
figclass: margin-caption
alt: My figure text
name: particle_collision_blank
---
MARGINCAP
\end{verbatim}

Let's add it up:

Let's imagine a chart on a whiteboard. Here is the total mass-energy of
the two beam protons:

\begin{verbatim}
---
width: 90%
figclass: margin-caption
alt: My figure text
name: justmass
---
MARGINCAP
\end{verbatim}

The vertical axis will tally energies of each of the objects scaled in
units of our fake units of \(m_pc^2\):

\begin{enumerate}
\def\labelenumi{\arabic{enumi}.}
\item
  The mass-energies of each proton.
\item
  This totals to 2.
\item
  The total mass-energy required in order to make two things is
  \(2 \times 3.5 = 7\).
\end{enumerate}

Now let's add in the kinetic energies:

\begin{verbatim}
---
width: 90%
figclass: margin-caption
alt: My figure text
name: massandK
---
MARGINCAP
\end{verbatim}

\begin{enumerate}
\def\labelenumi{\arabic{enumi}.}
\item
  Each beam proton had \(K=3\) .
\item
  So the total of the beams' energies of motion is 6.
\end{enumerate}

Putting it all together:

\begin{verbatim}
---
width: 90%
figclass: margin-caption
alt: My figure text
name: massKandtotal
---
MARGINCAP
\end{verbatim}

\begin{enumerate}
\def\labelenumi{\arabic{enumi}.}
\item
  The kinetic energies plus the total mass energies of the protons
\item
  totals to 8.
\end{enumerate}

As you can see, we needed \(E_T = 7\) but the collision provided a total
energy of \(8\), more than required in order to just make two things.

\begin{verbatim}
:class: warning

**Glad you asked:** The things are not just hanging around after they're produced. There is energy left over after they are created as real particles...and so they're moving. They acquire some kinetic energy.
\end{verbatim}

The two, shiny new "things" have kinetic energies of their own of
\(K_\text{things} = 1\).

:::

What we just did is solve the energy conservation equations:

\begin{align*}
E_{T,0}(p, L) + E_{T,0}(p, R) =& E_{T}(\text{thing 1}, L) + E_{T}(\text{thing 2}, R) \\
E_T = 2 \times m_pc^2 + 2 \times K_p =& 2 \times m_tc^2 + 2 \times K_t \\
2 \times 1 + 2 \times 3 =& 2 \times 3.5 + 2 \times K_t \\
2 + 6 =& 7 + 2K_t \\
2K_t =& 8-7 = 1
\end{align*}

A lot of relativity is pretty elementary math.

That. In a nutshell is much of the goal of particle physics: make new
things by using relativity (and quantum mechanics) out of very fast old
things.

\begin{verbatim}
:class: danger

<a href="https://loncapa.msu.edu//tiny/msu/c53cQ2"  target="_blank">Making new things</a>
\end{verbatim}

\hypertarget{some-chemistry-as-viewed-by-a-physicist}{%
\subsubsection{Some Chemistry: as viewed by a
physicist}\label{some-chemistry-as-viewed-by-a-physicist}}

Remember the hydrogen atom is one proton and one electron and that
electron and proton are bound together by the electrostatic attraction
due to their opposite electronic charges. Maybe you remember from
chemistry that in order to liberate the electron -- to "ionize"
hydrogen, you must supply a particular minimum energy of 13.6 electron
volts or \(2.18 \times 10^{-18}\) J.

You would have learned that the electron is "bound" to the proton with a
potential energy that's negative. The energy diagram would look like
this in a chemistry class on the left and a physics class on the right:

\begin{verbatim}
---
width: 90%
figclass: margin-caption
alt: My figure text
name: hydrogenCP
---
MARGINCAP
\end{verbatim}

Let's look at each item here by the circled numbers in the diagram.

The left side, the chemist picture:

1. The hydrogen energy level diagram is a succession of energy levels
from the lowest -- which is the most stable -- to the highest, which is
complete ionization....liberation of the electron.

2. For chemists, the zero of energy is when the electron has been
liberated...that's zero chemical energy.

3. When the electron is bound, it has a negative electrostatic potential
energy of \(-13.6\) electron volts.

4. Give an electron a hit of \(+13.6\) eV (like from a photon), and it
will be promoted away from the proton and is no longer bound: it's a
separate proton and a separate electron.

The right side, the physicist picture:

5. For the physicist the zero of energy is where there is no motion and
there is no mass.

6. For hydrogen, the energy is all in the mass of the hydrogen
\underline{atom}.

7. For a real proton and a real electron, which are not bound together
into an atom, the energy is just \(m_pc^2 + m_ec^2\). Notice that this
value goes to the zero value of the chemist picture: an independent and
isolated proton and electron.

Why does the electron stay bound to the proton in hydrogen in the
relativistic picture? Because the mass of a hydrogen atom is less than
the mass of an electron and a proton. So it cannot just rattle apart
into those two particles since

\[E_m(\text{hydrogen atom}) < E_m(\text{real proton}) + E_m(\text{real electron})! \nonumber\]

So how to reconcile this? Remember that what's not in the picture is the
electric field.

\begin{itemize}
\item
  We know from Maxwell that the electric field has energy.
\item
  We know from Einstein that energy and mass are equivalent.
\item
  So the "missing mass" in the hydrogen atom is resident in the energy
  of the electric field.
\end{itemize}

It's all accounted for!

\begin{verbatim}
---
width: 90%
figclass: margin-caption
alt: My figure text
name: hydrogenbalance
---
MARGINCAP
\end{verbatim}

Since, after all, energy is mass and mass is energy. This same idea is
what happens inside of nuclei and that "mass deficit" can be liberated,
either in a controlled way -- nuclear power -- or an explosive way --
nuclear weapons.

Everything you've learned about energy can be cast in this light.
Potential energy can be done away with:

\begin{itemize}
\item
  Raise a mass above the ground? You increase its mass. That's what
  we've called potential energy.
\item
  Heat up your eggs in the morning? They become more massive. The
  jiggling of the molecules in the eggs which we would call heat, can be
  interpreted as an overall increase in the eggs' mass.
\end{itemize}

Energy suddenly becomes a simple idea! And, it's always conserved...as
long as mass-energy is included in the accounting.

\begin{verbatim}
:class: danger

<a href="https://loncapa.msu.edu//tiny/msu/f98qlQ"  target="_blank">Potential energies? Nah.</a>
\end{verbatim}

\hypertarget{isnt-anything-constant}{%
\subsection{Isn't Anything Constant??}\label{isnt-anything-constant}}

\begin{verbatim}
:class: warning

**Glad you asked:** So it would seem. But it's not quite like that and what's really important is actually what is always constant. That's where Einstein's frustrated university mathematics instructor comes in.
\end{verbatim}

As we learned from the introduction to this lesson, in 1908 Minkowski
put Special Relativity on a sophisticated mathematical foundation. His
formulation first frustrated Einstein who grumbled:

\begin{quote}
``{[}about{]} superfluous learnedness..." "Since the mathematicians have
grabbed hold of the theory of relativity, I myself no longer understand
it.''
\end{quote}

But he soon realized the importance of "invariance" in physics and made
use of that notion many times later.

\begin{verbatim}
:class: important

**Invariance**: something that is unchanged under some "transformation."
\end{verbatim}

I want to hint at some pretty mathematics from 50,000 feet. The idea of
a mathematical "transformation" is a very important idea and in fact
governs modern physics in a deep way.

Suppose we have a square in front of us. A perfect, mathematical square.
Let's pretend that we can mark the corners to keep track of them as we
play. But the numbers don't do anything but guide the eye. Each corner
is identical to the others.

A transformation is here a rotation of the square around its center, C.

\begin{verbatim}
---
width: 60%
figclass: margin-caption
alt: My figure text
name: square
---
MARGINCAP
\end{verbatim}

What kinds of transformations would leave the square "invariant"? That
is, if I ask you to close your eyes and then I rotate the square around
its center and ask you to open them, would you be able to tell whether
it changed or whether it stayed the same?

If I did a rotation like (a) above, you'd say that the square was not
the same. So we could conclude that a square is NOT invariant against a
transformation of +45º about its center.

If I did a rotation like (b) above, you would not be able to tell that
anything changed. Indeed, a square is invariant with respect to a
transformation of 90º about its center. And 180º and 270º, 0º, and 360º.
These are the only transformations that leave a square invariant.

Minkowski was an expert in the branch of mathematics called then,
"Invariant Theory" and now "Group Theory" and he understood well that
the manipulative description that I just gave could be constructed as
math equations -- functions -- for which a transformation on the
function leaves it unchanged in its form.

For example, remember that the equation for a circle is:

\[R^2 = x^2 + y^2 \nonumber\]

Suppose (in the same sense as you closing your eyes before the square
manipulation) we transformed the \(x\) values that define the edge of a
circle this way:

\[x \to -a \nonumber\]

Every \(x\) becomes a minus \(a\). What is the new equation?

\[R^2 = a^2 + y^2 \nonumber\]

This has the same FORM as the original. It doesn't matter what the
variables are called. In this Group Theory game, what matters is the
FORM of an equation. So a circle is invariant with respect to the
transformation of its variables into their negatives.

\hypertarget{lets-go-to-the-ballpark}{%
\subsubsection{Let's Go to the Ballpark}\label{lets-go-to-the-ballpark}}

Have you ever thought about how major league baseball stadiums are laid
out? No?

\begin{verbatim}
:class: warning

**Glad you asked:** Stay with me. Invariance is afoot.
\end{verbatim}

The MLB Rule Book says that...ahem:

\begin{quote}
Rule 1.04: "THE PLAYING FIELD: It is desirable that the line from home
base through the pitchers plate to second base shall run East
Northeast."
\end{quote}

Since baseball is not played in the morning, this orientation would
minimize the likelihood that a batter would be staring into the sun. The
advent of night baseball, indoor stadiums, and prevailing winds in
various locations led to modern era baseball diamonds to be oriented in
all sorts of different directions.

Here are two:

\begin{verbatim}
---
width: 90%
figclass: margin-caption
alt: My figure text
name: baseballorientation
---
MARGINCAP
\end{verbatim}

The left-hand figure is a picture of what the Rule Book demands and the
elderly Wrigley Field in Chicago was constructed precisely that way,
while the modern CoAmerica Park went its own way to accommodate its city
streets. (Again, batters will face east, but southeast. Still no sun.)

Is everything relative in baseball? Well, no! The distance from home
plate to the pitcher's rubber is 60 feet, 6 inches in every ball
park.\footnote{There are so many strikeouts in the current game (2021)
  that there's talk of taking some advantage away from the pitchers by
  lengthen the distance to the mound.}

Notice in the figure above that there are two \(x-y\) axes drawn on top
of each diamond, \(x_C-y_C\) for the Cubs and \(x_T-y_T\) for the
Tigers. Notice too that a circle has been drawn centered on home plate
-- the origin of the two coordinate systems -- in each of the parks.
That circle has a radius of 60' 6".

Let's write the formula that models that circle. There are two of them:

\begin{align*}
L_P^2(\text{Cubs}) =& x_C^2+y_C^2 \\
L_P^2(\text{Tigers}) =& x_T^2 + y_T^2  \\
\text{ but:  } L_P^2(\text{Cubs}) =& L_P^2(\text{Tigers})
\end{align*}

Clearly, these formulae look the same, and in the sprit of the above
discussion it doesn't matter what the names of the variables are, the
\emph{form} of the equations are identical. We would say that the length
is invariant with respect to the choice of coordinate system since the
\(L_P\) for each ballpark is the same 60' 6".

In a flat space of the sort you might have learned the geometry of in
high school, for any number of coordinate systems, we could write:

\[L^2 = x_1^2 + y_1^2 = x_2^2 + y_2^2 = x_3^2 + y_3^2 = ... \nonumber\]

where now the different coordinate systems are called 1, 2, or 3 and not
Cubs, Tigers, Yankees...

This is very much like a well-worn example: Suppose a town hires two
contractors to work on the creation of a set of streets in town. The
first contractor uses magnetic north for his surveying jobs and the
second contractor uses polar north for hers. The streets that they
proposed to construct would be identical, street to street, corner to
corner, house to house.

Let's look at such a circle as described by three different coordinate
systems.



  

That circle that's traced out when the axes are aligned is called the
Invariant Curve. You can see that an invariant length from your high
school geometry class is represented by a circle.

What's critical here is those plus signs. Any time that you find a
length to be invariant and modeled by:

\[L^2 = x_1^2 + y_1^2 = x_2^2 + y_2^2 = x_3^2 + y_3^2 = ... \nonumber\]

you know that you're working in the high school geometry -- called
Euclidean Geometry -- after the Greek mathematician who wrote the book
that everyone learned from for 1500 years. \emph{Euclidean Geometry is
the geometry that governs points, lines, and planes on a flat surface.}
This will become critical in Einstein's future in more than one
surprising way and we'll draw the contrast more than once as we move
through his work.

Here's where Einstein's frustrated math professor comes in. Let's have a
baby.

\hypertarget{minkowskis-geometry}{%
\subsubsection{Minkowski's Geometry}\label{minkowskis-geometry}}

What we've seen so far is that relativity is going to mix space and
time.

We call the relativistic combination of space and time: Spacetime. 

A natural question is what's the "length" in space and time:
\emph{Spacetime}? What's the invariant curve for relativity?

\begin{verbatim}
What would that mean? It would be a way to describe an invariant interval as viewed from all possible reference frames. Following around the invariant curve would take you from Home to Away to another Away and another.
\end{verbatim}

Let's try the Euclidean Invariant curve for the "length" of space and
time.

\begin{verbatim}
Space and time are similarly dealt with so far, but they have different dimensions, right? Space might be measured in meters and time might be measured in seconds. Not the same. But we have a way around that. By convention we'll use a "space" dimension for time by always multiplying t by c.
\end{verbatim}

\textbf{A proposal:} We might guess that the "length" of a space and
time interval (we use "\(\Delta s\)" to represent a spacetime interval)
that the space and time coordinates would lead to a circle in Spacetime
just like our "regular" space of ballparks and surveyors. For our "1"
and "2" relatively moving coordinate systems, we might try

\[\Delta s^2 (\text{proposal}) = (c\Delta t_1)^2 + \Delta x_1^2+ \Delta y_1^2 + \Delta z_1^2 = (c\Delta t_2)^2 + \Delta x_2^2+ \Delta y_2^2 + \Delta z_2^2. \nonumber\]

Seems reasonable, but let's go to the hospital and think about what a
Spacetime geometry of this sort might look like.

On the left of the figure below is then our proposal. (Notice that we
can easily distinguish the future, with times greater than zero, from
the past, where time would be negative given the choice of the origin.)

On the right is a "situation" for us to think about with that
assumption.

\begin{verbatim}
---
width: 90%
figclass: margin-caption
alt: My figure text
name: cry
---
MARGINCAP
\end{verbatim}

So we're in the delivery room at a hospital. At the \(x=0\) point the
newly born baby is first...well, born...and then some time later at that
same space point, she cries, at the point labeled \(A\). Normal
biological stuff. It's a cause for celebration and little while later in
time everyone in the family cries.

Suppose we're driving by the hospital and somehow (don't ask how) we
observe that same blessed event. Now as time progresses, so does our
distance from the origin, so we all agree that the baby is born at the
origin but we see the first cry at point \(B\), since we've moved in our
frame. That's maybe not too hard to imagine.

Now we drive by in a different direction -- a different frame, that
Special Relativity should be able to describe. We again see the baby
born at the origin but uh-oh. \emph{Before the birth, the baby cries}.
If the circle is the space and time invariant curve then the rules of
relativity should work for all points on the curve -- this proposed
circle.

But that's not how births happen! The first crying event for every
newborn is after, not ever before, the birth. \emph{We have to conclude
that the Euclidean circle is not the invariant curve for Special
Relativity.} In the abstract, this is not just a silly example, it
implies \emph{a violation of cause and effect.} Where an effect seems to
happen before its cause. Not good.

While the maternity ward was not on Hermann Minkowski's mind, he did
figure out the correct kind of geometry to describe events in Special
Relativity. It was a new geometry that he invented and today we call the
space in which that geometry functions, Minkowski Space.

\hypertarget{spacetime}{%
\subsection{Spacetime}\label{spacetime}}

Let's get familiar with Spacetime. My ability to draw in 4 dimensions is
nil. So we'll write only time horizontally (remember \(ct\)) and \(x\)
vertically, proudly representing all three space dimensions. Stand
proud, \(x\)!

\hypertarget{digging-into-spacetime}{%
\subsubsection{Digging Into Spacetime}\label{digging-into-spacetime}}

Look at this figure with time on the horizontal axis (remember,
\(c \times t\) and one dimension of space on the vertical.

\begin{verbatim}
---
width: 90%
figclass: margin-caption
alt: My figure text
name: spacetimebare
---
MARGINCAP
\end{verbatim}

You've seen this many times. The origin of this plot is where we are
\underline{right now}. All frames would agree on this point as we've
done with trains in our earlier lessons.

The 45º angles are important since the slope would be

\begin{align*}
\text{slope } =& \frac{\Delta x}{\Delta ct} \\
\text{for } \frac{\Delta x}{\Delta t} =& c, \\
\text{slope } =& 1
\end{align*}

So a trajectory in spacetime for an object that moves at the speed of
light is 1. This is a light-like, or "light curve." These diagonals
create a multidimensional volume that looks like a cone, the
\textbf{Light Cone}.

\begin{verbatim}
---
width: 90%
figclass: margin-caption
alt: My figure text
name: spacetimelight
---
MARGINCAP
\end{verbatim}

So, the diagonals delineate four different regions of spacetime:

\begin{itemize}
\item
  In the future, where \(t>0\):

  \begin{itemize}
  \item
    inside of the light curve are objects that leave from us at
    \emph{less than the speed of light.} You know, normal stuff. We'd
    call that region the \textbf{Absolute Future}.
  \item
    outside of the light curve are objects that leave from us at
    \emph{greater than the speed of light}. You know, stuff of science
    fiction! We'd call that region the \textbf{Absolute Nowhere}.
  \end{itemize}
\item
  In the past, where \(t<0\):

  \begin{itemize}
  \item
    inside of the light curve are objects that move at \emph{less than
    the speed of light.} You know, normal stuff. We'd call that region
    the \textbf{Absolute Past}.
  \item
    outside of the light curve are objects that move at \emph{greater
    than the speed of light}. You know, stuff of science fiction! We'd
    call that region the \textbf{Absolute Nowhere}.
  \end{itemize}
\end{itemize}

The trajectories of all objects follow their "\textbf{Worldline}" which
is of course inside of the light cone.

\begin{verbatim}
---
width: 90%
figclass: margin-caption
alt: My figure text
name: spacetimeworldline
---
MARGINCAP
\end{verbatim}


🖋 📓

Notice one other thing about this:

\begin{verbatim}
There are regions of spacetime at any point which are unreachable from the past – could not have influenced us – and unreachable in the future – we can not communicate with them. The universe has expanded fast enough by now that there are places in our universe with which we have no contact.
\end{verbatim}

It's inside this new space of spacetime that we now need to address the
lingering question of what is the Invariant Curve for spacetime? That's
where Einstein's frustrated university math professor enters the scene.

\hypertarget{minkowski-space}{%
\subsubsection{Minkowski Space}\label{minkowski-space}}

We know that the circle will not work in Spacetime. Otherwise, we'd be
seeing babies everywhere before they arrive.

What Minkowski rigorously found was that a Spacetime length is not
represented by

\begin{verbatim}
:class: warning

$$
\Delta s^2 (\text{proposal}) = (c\Delta t_1)^2 + \Delta x_1^2+ \Delta y_1^2 + \Delta z_1^2 = (c\Delta t_2)^2 + \Delta x_2^2+ \Delta y_2^2 + \Delta z_2^2. \nonumber
$$

but rather by:

$$
\begin{align*}
\Delta s^2 (\text{Spacetime}) = (c\Delta t_1)^2 - \Delta x_1^2- \Delta y_1^2 - \Delta z_1^2 = (c\Delta t_2)^2 - \Delta x_2^2 - \Delta y_2^2 - \Delta z_2^2. \nonumber
\end{align*}
$$

This is not the equation of a circle! See those **negative signs**? Let's make it just one dimension:

$$
\Delta s^2 (\text{Spacetime}) = (c\Delta t_H)^2 - \Delta x_H^2 \nonumber
$$
\end{verbatim}

The \(H\) reminds us that we're talking about a spacetime length in some
\(A\) frame moving frame as viewed by us, \(H\), maybe lots of us moving
relative to one another and that \(A\) frame. All inertial observers
would agree on the length of a spacetime interval according to this
formula, just like we would all agree on the length from home plate to
the pitcher's rubber in flat "regular" space.

\begin{verbatim}
This is the equation of a **hyperbola**. Spacetime is hyperbolic in shape.
\end{verbatim}

Let's plot this on our spacetime diagram.

\begin{verbatim}
---
width: 90%
figclass: margin-caption
alt: My figure text
name: spacetime0.5
---
MARGINCAP
\end{verbatim}

On the left in the figure is a situation for all observers of some event
that happened at a time \(0.5\) units of time in the Home frame. You can
see that in the blowup around the origin where for no change in \(x\)
(after all, it's in the "stationary" frame). Every point on that red
curve are spacetime points that can be reached for a "proper time" of
0.5. That's the time in the rest frame of the object that's moving
relative to all of the others. You can see that it's 0.5 as it's the
shortest distance between the origin and the curve and where there is no
change in the space dimension -- a horizontal line at \(ct=0.5\) .

Now let our object move a little further (live a little longer in the
Away frame), to \(ct=1\).

\begin{verbatim}
---
width: 90%
figclass: margin-caption
alt: My figure text
name: spacetime1.0
---
MARGINCAP
\end{verbatim}

Here's the counter-intuitive thing about an hyperbola. Let's go back to
the airport. The moving sidewalk has a constant speed, which we could
(and did) plot on our graph as a positively sloping line. We can draw
arrows on the diagram from our common point of when time starts to some
time later. For the weary traveler, that line appears to be longer than
it is for the couch people, but in that formula of the relativistic
interval, \(\Delta s^2\) are the same for both. We'll work that out in
an example.

So, there are three important quantities that are NOT relative in
Special Relativity!

\begin{itemize}
\item
  The speed of light, \(c\).
\item
  The interval, \(\Delta s^2 = (c\Delta t)^2 - (\Delta x)^2\).
\item
  The "invariant mass," \((\Delta mc^2)^2 = E_T^2 - (pc)^2\).
\end{itemize}

Even this silly equation expresses a common feature with the last two
above:

\[(\text{Bob})^2 = (\text{Mary})^2 - (\text{Sam})^2 \nonumber\]

It doesn't matter what the variables are called. That the form of the
equations are the same expresses a deep mathematical feature of that new
kind of geometry. That's what Minkowski found.

\hypertarget{the-aftermath}{%
\subsection{The Aftermath}\label{the-aftermath}}

It might be worth reviewing a little of Einstein's personal timeline
since 1905.

\begin{itemize}
\item
  1906: Promoted to Technical Expert Second Class.
\item
  1906: He published six papers.
\item
  1907: He published 10 papers.
\item
  1907: Establishes the Principle of Equivalence (stay tuned)
\item
  1908: Turned down as a high school teacher and completes his
  \emph{habilitation} thesis in order to qualify for entry university
  teaching.
\item
  1908 (Finally) appointed as a lecturer at the University of Bern
\item
  1908: Minkowski's geometrical interpretation of Spacetime.
\item
  1909: Resigns from the Patent Office(!)
\item
  1909: Appointed Associate Professor of Theoretical Physics at the
  University of Zurich
\end{itemize}

During the years leading up to 1905 -- five years -- he worked alone,
except for fruitful discussions with his Olympia Academy friends. He
complains about not having access to any university library, so he
cannot read scientific papers.

\hypertarget{the-real-physics-olympia-berlin}{%
\subsubsection{The Real Physics Olympia:
Berlin}\label{the-real-physics-olympia-berlin}}

\begin{quote}
Einstein's devoted sister wrote that Einstein had hoped for some
recognition: "...he was bitterly disappointed...icy silence followed the
publications."
\end{quote}

In the western world at the turn of the 20th century, there was no
institution that remotely compared with the University of Berlin. And
there was no higher authority in physics, especially in Theoretical
Physics, than the Professor of Theoretical Physics, Max Planck. I'll
have lots to say about Planck when we get to Quantum Theory, but suffice
it to say, if there was ever an 800 pound gorilla in Theoretical Physics
in 1905, it was Planck. That turns out to be important.

Planck was on the editorial board of the primary physics journal in the
world, \emph{Annalen der Physik} and was in charge of theoretical
submissions. So all of Einstein's 1905 papers required his approval.
This would be hard to do today.

Einstein didn't even qualify to be classified as prominent as
"unknown"...he could not have been lower on any professional rung. He
had no PhD. He had no scientific job. He suffered the opposite of good
recommendations.

And yet, Max Planck took his papers seriously and by 1906, was giving
seminars on Relativity and published a correct derivation of the
relativistic momentum relation. He and Albert began corresponding. His
"bitter disappointment" turned into enthusiasm. Planck publicly defended
a subtle point of criticism of Relativity and Einstein joined in, along
side of his senior correspondent. Planck replied in January of 1907
with:

\begin{quote}
``As long as the proponents of the principle of relativity constitute
such a modest little band as is now the case...it is doubly important
that they agree among themselves.''
\end{quote}

Planck suggested that he would like to travel to Bern to meet Einstein
in person. Planck could not make the trip and so he sent his assistant,
Max Laue. It was only upon preparing to go to Bern that Laue (and
Planck) discovered that Einstein was not at the University of Bern...nor
at any university at all. When he got to the third floor of the Post and
Telegraph building, "The young man who came to meet me made so
unexpected impression on my that I did not believe he could possibly be
the father of the relativity theory...so I let him pass." Einstein
wandered back through the reception area and they made contact.

They spent hours together and eventually Laue published eight papers
himself on relativity and they became fast friends.

\hypertarget{nothing-is-simple}{%
\subsubsection{Nothing Is Simple}\label{nothing-is-simple}}

The academic ladder in Europe at the time, and still in many places
today, requires not only an original PhD degree, but a second PhD called
the "\emph{habilitation}." Einstein didn't want to bother, but failed at
the high school job application that he made he succumbed to convention,
qualifying for the degree at Bern. Soon a more prominent position opened
at the University of Zurich where Alfred Kleiner persuaded the
university to create a position in Theoretical Physics, albeit at the
associate professor level. Two applied: Einstein, whom Kleiner knew and
liked, and Friedrich Adler a friend of Einstein's from their college
days.

The job was to go to Adler, but he actually went to Kleiner to complain
that it should instead go to Einstein. Einstein's reputation for having
zero people skills made that a difficult decision for Kleiner but he
went to Bern to hear Einstein teach a class. It was terrible. Einstein
knew it and regrettably word of it got out around Europe: Einstein was
"a long way from being a teacher." Einstein was furious and wrote to
Kleiner complaining that vicious rumors were being spread about him.
About Adler:

\begin{verbatim}
Adler was an enigma, and a murderer. His father was a powerful official in the  Austrian Social Democratic party, to which Friedrich was a zealot participant. In 1911 he gave up science to become a journalist and partisan. He was against the Austria-Hungary war policy and on October 21st in 1916 young Adler calmly entered a hotel restaurant and shot and killed the Austrian minister-president Count Karl von Stürgkh, three times with a pistol. He made no effort to flee and admited his guilt. He was sentenced to death, but that was commuted by Emporer Charles to 1-18 years in prison. He was released at the end of WWI and lived to be 80 years old in 1960.
\end{verbatim}

Kleiner acquiesced and advised Einstein to try one more time, this time
in Zurich. Einstein did it. "Contrary to my habit, I lectured well on
that occasion." In spite of faculty reservations about Einstein's
"Jewishness" he was offered the job, which paid less than his Patent
Office position. "So, now I too am an official member of the guild of
whores."

Finally, by 1908 - eight years after he completed college and three
years after he changed the world in multiple ways, Einstein was a fully
credentialed university professor.

\hypertarget{is-relativity-the-case}{%
\subsection{Is Relativity The Case?}\label{is-relativity-the-case}}

Of course, just being a novel idea is not sufficient for the acceptance
of a theory. It has to pass experimental tests and here too, nothing was
simple for Einstein.

Remember that Lorentz's theory of the ether was in many ways predictably
identical to Relativity. That included a Lorentzian view of mass which
also would predictably increase with speed (relative to the ether). In
1901 there were competing theories based on ideas of electromagnetism
that differed from Maxwell's and Lorentz'. Walter Kaufmann, whom we'll
learn more of later, took on the challenge to measure whether the mass
of an electron varied with velocity. More particularly, he measured the
so-called "charge to mass" ratio, which can be determined by
accelerating a charged particle through an appropriately designed
sequence of electric and magnetic fields. (We'll learn about this when
we discuss the discovery of the electron by the Englishman, J.J.
Thomson.)

Kaufmann was a very patient and skilled experimenter, but he got it
wrong. His measurements were analyzed by him to disfavor the Lorentz
prediction (and of course consequently, the 1905 Einstein prediction).
Kaufmann did experiments in 1901, 1902, 1903, and 1905 and concluded:

\begin{quote}
"\emph{The prevalent results decidedly speak against the correctness of
Lorentz's assumption as well as Einstein's. If on account of that one
considers this basic assumption refuted, then one would be forced to
consider it a failure to attempt to base the entire field of physics,
including electrodynamics and optics, upon the principle of relative
movement.} "
\end{quote}

In that same year, Planck took the highest velocity nine data points
from Kaufmann's paper and reanalyzed them and found multiple mistakes.
That led to many more measurements in various labs and by 1914 it was
clear: the Einstein (Lorentz) interpretation of the change of mass with
speed was confirmed and the alternatives were ruled out.

We've already talked about some of the classic tests of Special
Relativity and today, it's a tool and not "just" a theory. Its
combination with Quantum Mechanics has led to the most precise
predictions and confirmations of those predictions in history.

And yet, Einstein's Nobel Prize ignored relativity.

The mass measurements weren't the end of problems for Special
Relativity.

\hypertarget{jewish-science}{%
\subsubsection{"Jewish Science"}\label{jewish-science}}

Anti-semitism was in the fabric of European life

\begin{quote}
``As remarkable as Einstein's papers are\ldots it still seems to me that
something almost unhealthy lies in this unconstruable and impossible to
visualize dogma. An Englishman would hardly have given us this
theory\ldots{} the abstract conceptual character of the Semite expresses
itself.'' \textbf{Arnold Sommerfeld, Nobel Laureate}
\end{quote}

This was the beginning of the attack on ``Jewish Science'' by the
National Socialists Party in Germany..early 1930s

\begin{quote}
``Materialistic natural science had eclipsed the ``spiritual sciences,''
giving rise to the ``arrogant delusion`` that humankind can achieve the
``mastery of nature\ldots That influence has been strengthened by the
all-corrupting foreign spirit permeating physics and mathematics\ldots''
\textbf{Philipp Lenard, Nobel Laureate}
\end{quote}

\begin{quote}
``The scientist does not exist only for himself or even for his science.
Rather, in his work he must serve the nation first and foremost. For
these reasons, the leading scientific positions in the National
Socialist state are to be occupied not by elements alien to the Volk but
only by nationally conscious German men.'' \textbf{Johannes Stark, Nobel
Laureate}
\end{quote}

\begin{quote}
``His theory is not the keystone of a development, but a declaration of
total war, waged with the purpose of destroying what lies at the basis
of this development, namely, the world view of German man . . . This
theory could have blossomed and flourished nowhere else but in the soil
of Marxism, whose scientific expression it is, in a manner analogous to
that of cubism in the plastic arts and the unmelodies and unharmonic
atonality in the music of the last several years {[}``degenerate
science''!{]}. Thus, in its consequences the theory of relativity
appears to be less a scientific than a political problem.''
\textbf{Bruno Thüring}
\end{quote}

\begin{quote}
``The simple question that precedes every scientific enterprise is: who
is it who wants to know something, who is it who wants to orient himself
in the world around him? It follows necessarily that there can only be
the science of a particular type of humanity and of a particular age.
There is very likely a Nordic science, and a National Socialist science,
which are bound to be opposed to the Liberal--Jewish science, which,
indeed, is no longer fulfilling its function anywhere, but is in the
process of nullifying itself.'' \textbf{Adolf Hitler}
\end{quote}

Eventually, it became unsafe for Einstein to remain in Germany and he
emigrated to the United States in 1932.

{[}to do - the Opera experiment{]}

Let's go paint a house in Bern.

\hypertarget{what-to-remember-from-lesson-18}{%
\subsection{What To Remember from Lesson
18}\label{what-to-remember-from-lesson-18}}

\hypertarget{a-new-kind-of-energy}{%
\subsubsection{A New Kind of Energy}\label{a-new-kind-of-energy}}

Mass is a kind of energy. In an object's own rest frame, the value of
that energy is:

\[E_m = mc^2 \nonumber\]

As viewed by an observer in a co-moving inertial frame, that energy
would appear to be:

\[E_T=m\gamma c^2 \nonumber\]

This can be interpreted as observing that a moving object appears to
have more mass than its value measured within its own rest frame:

\[m_R = \gamma m \nonumber\]

where "relativistic mass" is sometimes used to represent that quantity.
In our terms, this might also be written:

\[m_H = \gamma m_A \nonumber\]

With some algebra, we can write the total energy of an object in another
way:

\[E_T =& E_m + K \\
K =& E_T - E_m \\
K =& m\gamma c^2 - mc^2 \\
K =& mc^2(\gamma -1)\]

where \(K\) is the relativistic kinetic energy.

\end{document}
